\subsection{Direct Preference Optimization}
\label{sec:dpo}

The RLHF pipeline described in Section~\ref{sec:rlhf} is effective but
operationally burdensome: it requires training a separate reward model,
maintaining both a policy and a reference policy in memory, and running an
online RL loop with PPO---each stage introducing its own hyperparameters and
instabilities.  \emph{Direct Preference Optimization}
(DPO)~\cite{rafailov2023dpo} eliminates the reward-modelling and
reinforcement-learning stages entirely by showing that the optimal policy under
the KL-constrained RLHF objective admits a closed-form solution, and that this
solution can be exploited to reparameterize the reward function in terms of the
policy itself.

%% ---------- Derivation ----------
\paragraph{Closed-form optimal policy.}
Recall that the KL-constrained RLHF objective seeks a policy $\pi$ that
maximises expected reward while remaining close to the supervised fine-tuned
reference policy $\pi_{\text{ref}}$:
\begin{equation}
\label{eq:rlhf-obj}
\max_{\pi}\;\mathbb{E}_{x \sim \mathcal{D},\, y \sim \pi(\cdot|x)}
  \bigl[r(x,y)\bigr]
  - \beta\, D_{\mathrm{KL}}\!\bigl(\pi(\cdot|x) \,\|\, \pi_{\text{ref}}(\cdot|x)\bigr),
\end{equation}
where $\beta > 0$ controls the strength of the KL penalty.  By writing the KL
divergence in its definition and taking the functional derivative with respect
to $\pi(y|x)$, one can verify that the unique optimum of
Eq.~\eqref{eq:rlhf-obj} is
\begin{equation}
\label{eq:optimal-policy}
\pi^{*}(y|x)
  = \frac{1}{Z(x)}\,\pi_{\text{ref}}(y|x)\,
    \exp\!\Bigl(\frac{1}{\beta}\,r(x,y)\Bigr),
\qquad
Z(x) = \sum_{y} \pi_{\text{ref}}(y|x)\,
       \exp\!\Bigl(\frac{1}{\beta}\,r(x,y)\Bigr).
\end{equation}
The partition function $Z(x)$ is intractable to compute for large vocabulary
spaces, but---crucially---it depends only on the prompt $x$ and cancels in the
key step that follows.

\paragraph{Reward reparameterization.}
Taking the logarithm of both sides of Eq.~\eqref{eq:optimal-policy} and
rearranging yields an expression for the reward in terms of the optimal policy:
\begin{equation}
\label{eq:reward-reparam}
r(x,y) = \beta\,\log\frac{\pi^{*}(y|x)}{\pi_{\text{ref}}(y|x)}
        + \beta\,\log Z(x).
\end{equation}
The reward is thus determined (up to the prompt-dependent constant
$\beta\log Z(x)$) by the log-ratio of the optimal policy to the reference
policy.  This reparameterization is the central insight of DPO: rather than
first learning a reward model and then optimizing a policy against it, one can
directly learn the policy that \emph{implicitly} defines the reward.

\paragraph{The DPO loss.}
Substituting Eq.~\eqref{eq:reward-reparam} into the Bradley--Terry preference
model~\cite{bradley1952rank}
$p(y_w \succ y_l | x) = \sigma\!\bigl(r(x,y_w) - r(x,y_l)\bigr)$,
the partition functions cancel, giving a probability of preferring the winning
completion $y_w$ over the losing completion $y_l$ that depends only on
policy log-ratios.  Maximising the log-likelihood of the observed preference
data $\mathcal{D} = \{(x^{(i)}, y_w^{(i)}, y_l^{(i)})\}$ under this model
yields the DPO objective:
\begin{equation}
\label{eq:dpo-loss}
\boxed{
\mathcal{L}_{\text{DPO}}(\theta)
  = -\,\mathbb{E}_{(x,\,y_w,\,y_l)\sim\mathcal{D}}\!
    \left[\log\sigma\!\left(
      \beta\left(
        \log\frac{\pi_{\theta}(y_w|x)}{\pi_{\text{ref}}(y_w|x)}
      - \log\frac{\pi_{\theta}(y_l|x)}{\pi_{\text{ref}}(y_l|x)}
      \right)
    \right)\right].}
\end{equation}
Training proceeds by standard gradient descent on the language-model parameters
$\theta$; no reward model, no value function, and no online sampling are
required.  The reference policy $\pi_{\text{ref}}$ is kept frozen throughout
and is typically the SFT checkpoint from which $\pi_{\theta}$ is initialised.

\paragraph{Practical advantages.}
DPO enjoys several compelling benefits relative to the full PPO-based RLHF
pipeline:
\begin{enumerate}[leftmargin=1.8em,itemsep=2pt]
\item \textbf{Architectural simplicity.}  The entire alignment procedure
      reduces to supervised learning on preference pairs.  There is no need to
      train or serve a separate reward model, and no RL optimiser (with its
      attendant clipping, GAE, and multiple rollout-to-update ratios) is
      required.
\item \textbf{Memory and compute efficiency.}  Because only the policy and
      reference model are loaded (and the reference can share weights via
      LoRA offsets), DPO typically requires 40--60\% less GPU memory and
      wall-clock time than PPO at comparable model scale.
\item \textbf{Training stability.}  The loss in Eq.~\eqref{eq:dpo-loss} is a
      simple cross-entropy--style objective; it does not suffer from the
      high-variance gradient estimates and reward hacking failure modes that
      plague online RL.
\end{enumerate}

\paragraph{Limitations.}
Despite its elegance, DPO has several well-documented shortcomings:
\begin{itemize}[leftmargin=1.8em,itemsep=2pt]
\item \textbf{Offline, static preference data.}  DPO optimises against a fixed
      dataset of preference pairs.  As $\pi_{\theta}$ drifts from the
      distribution that generated the data, the implicit reward signal becomes
      stale---a form of \emph{distribution shift} that degrades alignment
      quality over continued training.
\item \textbf{No online exploration.}  Because DPO never samples from the
      current policy during training, it cannot discover novel high-reward
      behaviours that lie outside the support of the static dataset.
\item \textbf{Likelihood displacement.}  Empirical studies have shown that DPO
      can decrease the absolute likelihood of \emph{preferred} completions
      while increasing the margin over dispreferred ones---a phenomenon termed
      \emph{likelihood displacement}.  This can degrade generation quality even
      as the preference-classification accuracy improves.
\item \textbf{Weaker long-horizon alignment.}  For tasks that require
      multi-step reasoning chains---where the correctness of later steps
      depends on earlier ones---the pairwise, completion-level preference
      signal in DPO provides limited credit assignment, making it less
      effective than methods that can assign step-level or
      outcome-level rewards.
\end{itemize}

\paragraph{Why DPO is insufficient for mathematical reasoning.}
The limitations above are particularly acute in the domain of mathematical
problem solving.  First, mathematical reasoning is \emph{verifiable}: a
derivation either reaches the correct answer or it does not, and this binary
outcome signal is far richer and cheaper to obtain than human preference
annotations.  DPO cannot exploit this signal because it requires pre-collected
\emph{pairwise} preferences rather than scalar correctness labels.  Second,
the solution space for mathematical problems is vast and compositional; a
model must explore diverse multi-step reasoning paths to improve, yet DPO's
offline nature precludes such exploration.  Third, the credit-assignment
problem is severe: a single algebraic error in the fifth step of a ten-step
proof renders the entire completion incorrect, but DPO's completion-level loss
provides no mechanism to localise fault.

These shortcomings motivate the search for alignment methods that (i)~leverage
verifiable outcome rewards rather than pairwise preferences, (ii)~sample from
the current policy during training to enable on-policy exploration, and
(iii)~achieve the computational simplicity of DPO without sacrificing the
performance benefits of online RL.  \emph{Group Relative Policy Optimization}
(GRPO), introduced in Section~\ref{sec:grpo}, addresses precisely this gap.
